\chapter{Portfolio Theory and Estimation}

\begin{objectives}
Derive efficient portfolios, understand estimation error, apply shrinkage and robust
construction, and evaluate portfolios through implementable out-of-sample decisions.
\end{objectives}

\section{Returns and wealth}

For simple return $R_t=(P_t+D_t)/P_{t-1}-1$, wealth compounds as
$W_T=W_0\prod_t(1+R_t)$. Log return $r_t=\log(1+R_t)$ adds through time.
Cross-sectionally, portfolio simple return is $R_{p,t}=w_{t-1}^\transpose R_t$ for
pre-return weights. Portfolio log return is not the weighted average of asset log
returns.

Annual arithmetic mean is often scaled by periods; volatility by square root of
periods under stationarity and suitable dependence. Long-horizon quantiles and
drawdowns do not obey a universal square-root rule.

\section{Mean--variance optimization}

For expected return $\mu$ and covariance $\Sigma$,
\[
\max_w\left\{w^\transpose\mu-\frac{\gamma}{2}w^\transpose\Sigma w\right\}
\quad\text{s.t.}\quad\ones^\transpose w=1.
\]
With a risk-free asset and unconstrained risky weights, the tangency direction is
$\Sigma^{-1}(\mu-r_f\ones)$. Its maximum squared Sharpe ratio is
\[
SR_{\max}^2=(\mu-r_f\ones)^\transpose
\Sigma^{-1}(\mu-r_f\ones).
\]

This is an ex ante population statement. Replacing inputs by noisy estimates can
produce extreme weights and poor realized performance.

\section{Estimation error}

Expected returns are much harder to estimate than covariance at ordinary horizons.
Optimization maps small errors through $\Sigma^{-1}$, emphasizing directions of
apparently low risk that often correspond to noise.

Sampling uncertainty is only one issue. Structural change, survivorship,
asynchronous markets, stale holdings, and benchmark revisions create model and data
uncertainty. Report weights and risk contributions, not just the objective value.

\section{Covariance shrinkage}

A linear shrinkage estimator is
\[
\widehat\Sigma_\delta=
\delta F+(1-\delta)S,\qquad \delta\in[0,1],
\]
where $S$ is sample covariance and $F$ a structured target. Shrinkage trades some
bias for lower variance and improved conditioning.

Factor covariance uses
$B\Sigma_fB^\transpose+D$. Robust covariance estimators reduce outlier sensitivity.
Exponentially weighted estimates adapt to changing volatility but reduce effective
sample size.

\section{Expected-return shrinkage and views}

Bayesian or empirical-Bayes estimation shrinks noisy asset means toward a grand mean
or factor-implied prior. Black--Litterman starts from equilibrium-implied excess
returns
\[
\pi=\delta\Sigma w_{\mathrm{mkt}}
\]
and combines views $P\mu=q+\text{noise}$ with a prior. Posterior mean depends on view
uncertainty and prior scaling; these are economic inputs, not tuning after seeing
performance.

\section{Constraints and regularization}

Long-only, leverage, box, sector, factor, turnover, and tracking-error constraints
stabilize portfolios while encoding mandate. A generic program is
\[
\min_w\frac12w^\transpose\widehat\Sigma w
-\lambda w^\transpose\widehat\mu
+\eta\norm{w-w_{\mathrm{prev}}}_1
\]
subject to exposure constraints. The turnover penalty approximates linear costs.
Constraints create shadow prices that quantify the value of relaxing a mandate under
the model.

\section{Resampling and robust portfolios}

Resampled efficiency averages optimized weights across simulated inputs. Distributionally
robust optimization protects against laws in an ambiguity set. Worst-case methods
can be conservative; uncertainty sets must be calibrated before observing desired
weights.

The one-over-$N$ portfolio is a serious benchmark. A complex method must earn its
estimation and trading costs relative to equal weight, global minimum variance, and
a simple factor-balanced allocation.

\section{Problems}

\begin{exercise}[Core: tangency]
Derive the tangency direction and maximum Sharpe ratio using Cauchy--Schwarz.
\end{exercise}
\begin{exercise}[Core: shrinkage]
Show that a convex combination of positive semidefinite covariance matrices remains
positive semidefinite. When is it positive definite?
\end{exercise}
\begin{exercise}[Advanced: sensitivity]
Differentiate unconstrained mean--variance weights with respect to $\mu$ and
$\Sigma$. Identify why small eigenvalues amplify error.
\end{exercise}


\chapter{Risk Measures, Backtesting, and Stress}

\begin{objectives}
Define coherent and convex risk measures, estimate VaR and Expected Shortfall,
backtest forecasts, and design scenarios that expose rather than conceal model
limitations.
\end{objectives}

\section{Loss conventions}

Let $L=-\Delta V$ be positive when money is lost. A risk number without horizon,
confidence, currency, valuation time, data window, and treatment of positions is
ambiguous.

\begin{definition}[Value at Risk]
\[
\VaR_\alpha(L)=\inf\{\ell:\Prob(L\leq\ell)\geq\alpha\}.
\]
\end{definition}

\begin{definition}[Expected Shortfall]
A general definition is
\[
\ES_\alpha(L)=
\frac{1}{1-\alpha}\int_\alpha^1\VaR_u(L)\dd u.
\]
For a continuous distribution, this equals
$\E[L\mid L\geq\VaR_\alpha(L)]$.
\end{definition}

\section{Coherent risk}

A risk measure $\rho$ is coherent when it is monotone, translation invariant,
positively homogeneous, and subadditive. Convex risk measures replace homogeneity
and subadditivity by convexity, allowing liquidity or concentration penalties.

\begin{theorem}[Expected Shortfall is coherent]
For integrable losses, $\ES_\alpha$ satisfies the four coherence axioms.
\end{theorem}
\begin{proof}
Use the representation
\[
\ES_\alpha(L)=
\inf_{z\in\R}\left\{
z+\frac{1}{1-\alpha}\E[(L-z)^+]
\right\}.
\]
Monotonicity and translation follow directly. Positive homogeneity follows by
scaling $z$. Subadditivity uses
$(X+Y-z_X-z_Y)^+\leq(X-z_X)^++(Y-z_Y)^+$ and takes infima.
\end{proof}

VaR is not generally subadditive for discontinuous or heavy-tailed losses. Under
elliptical distributions it may behave well, but coherence is a property over a
declared domain.

\section{Parametric and historical estimation}

Under normal $L\sim N(\mu_L,\sigma_L^2)$,
\[
\VaR_\alpha=\mu_L+\sigma_Lz_\alpha,\qquad
\ES_\alpha=\mu_L+
\sigma_L\frac{\phi(z_\alpha)}{1-\alpha}.
\]
Student tails, mixtures, filtered historical simulation, Monte Carlo, and EVT change
tail assumptions. Historical simulation places zero probability beyond observed
losses and assumes the selected past is representative.

Filtered historical simulation standardizes returns by estimated conditional
volatility, resamples residuals, then applies current volatility. This separates
scale dynamics from residual shape but inherits both model and resampling assumptions.

\section{Marginal and component risk}

For differentiable homogeneous $\rho(w)$, Euler's theorem gives
\[
\rho(w)=\sum_iw_i\frac{\partial\rho}{\partial w_i}.
\]
Component risk is $CR_i=w_i\partial\rho/\partial w_i$. It allocates the current
portfolio's risk but is not the same as incremental risk from deleting a position,
which is a finite change.

\section{VaR and ES backtesting}

For forecast $\widehat q_t$, exception $I_t=\ind_{\{L_t>\widehat q_t\}}$ should have
rate $1-\alpha$. Kupiec tests unconditional rate; Christoffersen tests transition
dependence. Quantile loss
\[
\ell_\alpha(q,L)=(\alpha-\ind_{\{L\leq q\}})(L-q)
\]
is strictly consistent for the quantile.

ES alone is not elicitable by a scalar loss, but $(\VaR,\ES)$ is jointly elicitable
under conditions. Compare models through proper scoring, calibration, and economic
consequences, not only exception counts.

\section{Stress testing}

Historical stress replays observed shocks. Hypothetical stress imposes coherent
factor moves. Reverse stress asks which scenario breaches a solvency or liquidity
threshold:
\[
\min_{\Delta x}\ \Delta x^\transpose\Omega^{-1}\Delta x
\quad\text{s.t.}\quad
\Delta V(\Delta x)\leq-L^\star.
\]
This finds a most plausible damaging shock under metric $\Omega$, not the only
dangerous scenario.

Scenarios should propagate rates, volatility surfaces, spreads, FX, liquidity,
defaults, collateral, and management actions consistently. Holding correlations
fixed in stress can be the main error.

\section{Liquidity risk}

Market liquidity concerns executable price and depth; funding liquidity concerns
ability to finance positions and margin. They interact through margin spirals.
Liquidity-adjusted loss includes liquidation horizon and impact:
\[
L_{\mathrm{liq}}=L_{\mathrm{market}}+
\sum_i\left(\frac{s_i}{2}Q_i+\eta_iQ_i^{3/2}\right).
\]
The horizon is position- and state-dependent. A ten-day scalar multiplier is not a
substitute for a liquidation schedule.

\section{Problems}

\begin{exercise}[Core: ES]
Derive normal Expected Shortfall by integrating the normal tail.
\end{exercise}
\begin{exercise}[Core: coherence]
Give a two-position discrete example where VaR violates subadditivity.
\end{exercise}
\begin{exercise}[Advanced: reverse stress]
Linearize portfolio loss and solve the Mahalanobis reverse-stress problem in closed
form. Then discuss gamma effects.
\end{exercise}


\chapter{Risk Allocation, Drawdowns, and Alternative Objectives}

\begin{objectives}
Construct risk-parity and drawdown-aware portfolios, distinguish volatility from
loss, and connect portfolio objectives to investor liabilities and horizons.
\end{objectives}

\section{Risk parity}

For volatility $\sigma(w)=\sqrt{w^\transpose\Sigma w}$, component contribution is
\[
RC_i=w_i\frac{(\Sigma w)_i}{\sigma(w)}.
\]
Equal-risk-contribution portfolios solve $RC_i=\sigma/n$ under positive weights.
One convex formulation minimizes
\[
\frac12w^\transpose\Sigma w-\sum_i b_i\log w_i
\]
and rescales the solution, where budgets $b_i>0$.

Risk parity does not make assets equally risky in all states. Covariance estimation,
volatility scaling, leverage, and concentration in a common macro factor determine
economic exposure. Levered bonds can dominate liquidity and inflation stress.

\section{Volatility targeting}

A strategy with raw return $r_t$ and forecast volatility $\widehat\sigma_t$ uses
exposure
\[
\ell_t=\min\left(\ell_{\max},
\frac{\sigma^\star}{\widehat\sigma_t}\right).
\]
It reduces exposure after volatility rises. Performance can improve through risk
stabilization and a negative return--volatility relation, but rapid deleveraging can
sell into stress and contribute to crowding.

The volatility estimate must use data through the decision time. Lag, rebalance
frequency, cap, financing, and turnover define the strategy.

\section{Drawdown}

For wealth $W_t$, running peak $M_t=\max_{s\leq t}W_s$ and drawdown
$D_t=1-W_t/M_t$. Maximum drawdown is path-dependent and its sample distribution
depends strongly on horizon. It is not a coherent one-period risk measure.

Conditional drawdown at risk and average drawdown measures focus on the distribution
of drawdowns. Drawdown constraints create dynamic policies because the state includes
the running maximum.

\section{Kelly and risk of ruin}

Log-optimal growth maximizes $\E[\log(W_{t+1}/W_t)]$. For repeated independent bets,
it maximizes asymptotic growth under its model. It can accept deep interim drawdowns
and is extremely sensitive to probability error and omitted tail outcomes.

Fractional Kelly multiplies estimated optimal exposure by $0<f<1$. This sacrifices
some modeled growth for reduced variance and parameter sensitivity. Hard ruin,
margin, or drawdown constraints should be modeled explicitly rather than expected to
emerge from log utility.

\section{Liability-driven investment}

An investor with liabilities cares about surplus
$S_t=A_t-L_t$. Asset-only volatility can be misleading: a long-duration asset may be
risky standalone but hedge a long-duration liability. Objectives include surplus
variance, funding-ratio shortfall, and expected utility of contributions.

Liability discounting itself can be regulatory, accounting, market-consistent, or
actuarial. The chosen curve changes both measured funding and hedge.

\section{Tracking error and active risk}

For benchmark $w_b$, active weights $a=w-w_b$ have ex ante tracking error
\[
TE=\sqrt{a^\transpose\Sigma a}.
\]
Information ratio is expected active return divided by tracking error. Fundamental
law approximations relate it to skill, breadth, and transfer coefficient under strong
independence and stationarity assumptions.

Constraints reduce the transfer from signal to portfolio. Active share measures
weight deviation but ignores covariance. Report both holdings-based and risk-based
measures.

\section{Problems}

\begin{exercise}[Core: ERC]
Derive first-order conditions of the log-barrier risk-budgeting objective and show
that risk contributions are proportional to $b_i$ after scaling.
\end{exercise}
\begin{exercise}[Core: surplus]
Derive the minimum-variance asset hedge of a random liability using covariance
projection.
\end{exercise}
\begin{exercise}[Advanced: drawdown state]
Formulate a dynamic program for terminal utility with a maximum-drawdown constraint.
Specify all state variables.
\end{exercise}


\chapter{Backtesting and Strategy Research}

\begin{objectives}
Build a point-in-time backtest, avoid data leakage, quantify costs and uncertainty,
and judge a strategy as a research claim rather than a chart.
\end{objectives}

\section{The event-driven timeline}

At each decision time:

\begin{enumerate}
\item ingest only records whose publication timestamp has passed;
\item update features using training-only transformations;
\item estimate or update the model;
\item freeze the target portfolio;
\item simulate orders in a stated execution window;
\item apply fills, costs, financing, borrow, and corporate actions;
\item mark positions and store an immutable audit row.
\end{enumerate}

Vectorized backtests are fast but can conceal timing. An event-driven ledger makes
cash, holdings, orders, fills, and value reconcile.

\section{Biases}

Look-ahead bias uses future information directly or through centered filters, revised
data, global standardization, or same-bar execution. Survivorship bias excludes
failed securities. Selection bias chooses favorable markets or periods. Data-snooping
bias chooses a model after examining outcomes. Publication bias hides failures.

Point-in-time universes, delisting returns, publication lags, and a research registry
are defenses. No statistical correction fully reconstructs undocumented exploration.

\section{Performance measurement}

For periodic excess returns $r_t^e$, estimated Sharpe ratio is
\[
\widehat{SR}=\frac{\bar r^e}{s_r}
\]
times an annualization factor under a stated frequency. Serial correlation changes
its standard error and invalidates naive square-root scaling of information.

Other metrics include Sortino ratio, maximum drawdown, Calmar ratio, hit rate,
skewness, tail ratio, turnover, capacity, beta, factor alpha, and certainty equivalent.
Choose predeclared metrics aligned with the strategy.

\section{Sharpe uncertainty}

Under iid normal returns, approximate standard error depends on sample size and
Sharpe itself. Non-normality and serial dependence require robust methods. A
probabilistic Sharpe ratio compares an observed estimate to a benchmark using
estimated skew and kurtosis; a deflated Sharpe also adjusts for selection among
trials.

Block bootstrap resamples consecutive observations to preserve local dependence.
Stationary bootstrap uses random block lengths. Strategy re-estimation should occur
inside each bootstrap if model-fitting uncertainty is part of the question.

\section{Costs and financing}

Net return includes gross asset P\&L, spread, commission, impact, borrow fees,
dividends on shorts, futures variation margin, funding, FX conversion, and tax if in
scope. Apply costs to trades, not average holdings. Use bid and ask where available.

Sensitivity should vary cost multipliers, execution delay, participation caps, and
assets under management. A strategy whose conclusion reverses under a small plausible
cost change is not robust.

\section{Benchmarking and attribution}

Benchmarks include cash, passive market, equal weight, factor-matched portfolios,
and the previous production model. Attribution can separate allocation, selection,
timing, factors, carry, convexity, and residual. Attribution is model-dependent and
must reconcile algebraically to total return.

\section{Decision protocol}

Before looking at the final test:

\begin{itemize}
\item state hypotheses and directional predictions;
\item freeze universe, features, estimator, tuning protocol, and costs;
\item define rejection and economic materiality thresholds;
\item name stress tests and subgroup analyses;
\item preserve code and configuration hash.
\end{itemize}

After the test, report every predeclared result and label new investigations
exploratory. A negative result with a clear mechanism test is valuable research.

\section{Problems}

\begin{exercise}[Core: timing audit]
Draw the complete timeline for a monthly accounting signal traded at the next
month's open. Include fiscal period end, filing, vendor ingestion, signal, and fill.
\end{exercise}
\begin{exercise}[Core: ledger]
Write identities reconciling starting cash, trades, fees, distributions, financing,
and ending portfolio value.
\end{exercise}
\begin{exercise}[Research]
Reproduce a published factor using point-in-time inputs. Compare gross, quoted-spread,
and impact-adjusted results, then test on an untouched geography.
\end{exercise}
